\documentclass[11pt,a4paper]{article}

\usepackage[margin=1in]{geometry}
\usepackage[T1]{fontenc}
\usepackage{lmodern}
\usepackage{microtype}
\usepackage{enumitem}
\usepackage{booktabs}
\usepackage[colorlinks=true,linkcolor=black,urlcolor=blue,citecolor=blue]{hyperref}
\usepackage{xurl}
\usepackage{titlesec}

\titleformat{\section}{\large\bfseries}{\thesection.}{0.6em}{}
\setlist{itemsep=2pt,parsep=0pt,topsep=4pt}

\title{\textbf{Weekly Progress Report}\\[0.3em]
\large Weekend of Saturday 18  Sunday 19 July 2026}
\author{Ahmed Khader}
\date{19 July 2026}

\begin{document}

\maketitle

\begin{abstract}
\noindent This report covers the work carried out over Saturday and Sunday.
The weekend was spent closing a knowledge gap identified during the Thursday
sessions: I mapped out the technical and mathematical background required,
built a structured study plan around it, and completed Phase~1 of that plan.
Phases~2 and~3 are scheduled for Monday, after which paper review begins.
\end{abstract}

\section{Background: What initiated This Work}

After attending the Thursday sessions, it became clear that parts of the
discussion sat on foundations I had not yet covered. Rather than continue
forward with partial understanding, I decided to stop and address the gap
directly.

\section{Building the Plan}

I began by researching the topic in depth, gathering material from across the
internet using both conventional search and AI models to identify what a
complete coverage of the subject would actually require, including the
equivalent mathematical background that underpins it.

I then fed the collected material to Claude and worked with it to turn a loose
pile of sources into a structured, phased learning plan. The result is a plan I
can walk through step by step rather than reading at random.

The wider aim is to build the background needed to read papers properly: to
follow them at depth, to judge what contribution each one actually makes, and
to appreciate the complexity of the problems they tackle. I also expect this
foundation, and the mathematics that comes with it, to pay off later when it
comes to writing paper.

\begin{itemize}
  \item \textbf{Study plan (Claude conversation):}\\
        \url{https://claude.ai/share/2c51b57a-d22f-4987-a8cc-5614a9044228}
  \item \textbf{Supporting reference:} Harvard notes covering one of the core
        topics:\\
        \url{https://people.math.harvard.edu/~knill/teaching/math118/118_dynamicalsystems.pdf}
\end{itemize}

\section{Progress to Date}

\begin{itemize}
  \item \textbf{Phase 1: Complete.} The foundational material identified in
        the plan has been worked through.
  \item \textbf{Phases 2 and 3: In progress.} Both are scheduled to be
        completed on Monday.
\end{itemize}

\section{Plan for the Coming Week}

With the foundations in place, the focus shifts from acquiring background
knowledge to applying it.

\begin{itemize}
  \item \textbf{Monday:} Finish Phases~2 and~3 of the study plan.
  \item \textbf{Tuesday:} Begin gathering papers and scanning them in depth,
        producing feedback on each.
  \item \textbf{Tuesday/Wednesday onward:} Assigned tasks proceed in parallel
        with the paper review.
  \item \textbf{By end of week:} All required background knowledge acquired;
        20+ papers scanned, summarised, and reviewed.
\end{itemize}

This section will be updated in every report sent, so that it always reflects
the work actually completed against the plan.

\end{document}
